\documentclass[11pt]{amsart}
\usepackage[margin=1.1in]{geometry}
\usepackage{amsmath,amssymb,amsthm,booktabs,hyperref,enumitem}
\hypersetup{colorlinks=true,linkcolor=blue,citecolor=blue,urlcolor=blue}

\newtheorem{theorem}{Theorem}
\newtheorem{lemma}[theorem]{Lemma}
\newtheorem{proposition}[theorem]{Proposition}
\newtheorem{corollary}[theorem]{Corollary}
\theoremstyle{definition}
\newtheorem{definition}[theorem]{Definition}
\theoremstyle{remark}
\newtheorem{remark}[theorem]{Remark}

\newcommand{\N}{\mathbb{N}}
\newcommand{\sqf}{squarefree}
\DeclareMathOperator{\lcm}{lcm}

\title[Erd\H{o}s Problem \#848 for all $N$]{Erd\H{o}s Problem \#848 holds for every $N$:\\ a computer-assisted completion of Sawhney's theorem}
\author{Ben Zanghi}
\date{Draft of 12 October 2026 --- not yet refereed}

\begin{document}

\begin{abstract}
Erd\H{o}s and S\'ark\"ozy asked whether the largest $A\subseteq\{1,\dots,N\}$ such that $ab+1$ is never
squarefree for $a,b\in A$ is given by the residue class $7 \bmod 25$. Sawhney proved this for all sufficiently
large $N$; the best explicit threshold previously available was $N_0=\exp(1420)$. We prove the statement for
every $N\ge 1$ (with $a=b$ allowed). The proof has three parts: an exact computation for $N\le 10^4$, a
dynamically maintained colouring certificate for $10^4\le N\le 10^9$ that is replayed by an independent
verifier, and a new elementary explicit argument for $N\ge 5\cdot 10^8$ whose numerical inequalities are
checked in interval arithmetic. The argument for large $N$ is self-contained and does not use Sawhney's
stability theorem.
\end{abstract}

\maketitle

\noindent\textbf{Disclosure.} This work was carried out with substantial assistance from an AI system
(Claude, Anthropic), which wrote the code and drafted the argument and this manuscript. All computations
are reproducible from the accompanying code; see Section~\ref{sec:repro}. An independent AI review
re-implemented the analytic bounds without access to our code. The manuscript has not yet been checked by a
human expert, and the author welcomes scrutiny.

\section{Introduction}

For $N\ge1$ call $A\subseteq[1,N]$ \emph{admissible} if $ab+1$ is not squarefree for all $a,b\in A$
(including $a=b$). Let
\[
  M(N)=\#\{n\le N:\ n\equiv 7 \pmod{25}\}.
\]
The set $P_7(N)=\{n\le N: n\equiv 7\ (25)\}$ is admissible since $ab+1\equiv 50\equiv 0\pmod{25}$.
Erd\H{o}s and S\'ark\"ozy \cite[p.~239]{Er92b} asked whether it is maximal; this is Problem~\#848 in Bloom's
database \cite{EP848}. Van Doorn observed that every admissible set has density at most
$2\sum_{p\equiv1(4)}p^{-2}\approx 0.108$, Weisenberg improved this to $\approx0.105$, and Sawhney
\cite{Sa25} proved that the answer is yes for all sufficiently large $N$, together with a stability
statement. An explicit version of Sawhney's argument, produced with AI assistance by Sothanaphan and others,
gives the threshold $N_0=\exp(1420)$ \cite{So26}. Direct computation had been carried out for $N$ up to a few
hundred.

\begin{theorem}\label{thm:main}
For every $N\ge1$ and every admissible $A\subseteq[1,N]$ we have $|A|\le M(N)$. In particular, the
residue class $7\bmod 25$ is a maximum admissible set for every $N$.
\end{theorem}

\begin{remark}
If one only requires $ab+1$ to be non-squarefree for $a\ne b$, the statement is different for small $N$:
Cambie noted that $\{19,33,35,46,53\}$ is such a set. Our argument uses the diagonal condition in an essential
way, through the restriction that $a^2+1$ is non-squarefree, and does not address that variant.
\end{remark}

\subsection*{Strategy}
Let $G_N$ be the graph whose vertices are the $n\le N$ with $n^2+1$ non-squarefree, with $a\sim b$
(for $a\ne b$) iff $ab+1$ is non-squarefree. Admissible sets are exactly the cliques of $G_N$, so
Theorem~\ref{thm:main} says $\omega(G_N)=M(N)$. We prove it in three overlapping ranges:
\begin{center}
\begin{tabular}{@{}lll@{}}
\toprule
range & method & section\\
\midrule
$1\le N\le 10^4$ & exact clique number (two independent programs) & \ref{sec:small}\\
$9982\le N\le 10^9$ & colouring certificate + independent verifier & \ref{sec:cert}\\
$N\ge 5\cdot 10^8$ & explicit inequalities (interval arithmetic) & \ref{sec:analytic}\\
\bottomrule
\end{tabular}
\end{center}

\section{Preliminaries}\label{sec:prelim}

Since $n^2\equiv 0,1\pmod 4$, we have $4\nmid n^2+1$. If $p\equiv3\pmod4$ then $p\nmid n^2+1$. Hence $n^2+1$ is
non-squarefree iff $q^2\mid n^2+1$ for some prime $q\equiv1\pmod4$. For such $q$ let
$S_q=\{n:\ q^2\mid n^2+1\}$, which is the union of two residue classes $\pm i_q\bmod q^2$. Since
$25\mid n^2+1\iff n\equiv\pm7\pmod{25}$, the vertex set of $G_N$ is the disjoint union
\[
  P_7=\{n\equiv7\},\qquad P_{18}=\{n\equiv18\},\qquad R\subseteq \bigcup_{q\ge13,\ q\equiv1(4)} S_q,
\]
where all three are intersected with $[1,N]$ and $R$ denotes the remaining vertices. Note
$|P_{18}|\le|P_7|=M(N)$, and both $P_7$ and $P_{18}$ are cliques.

For a clique $A$ we write $X=A\cap P_7$, $Y=A\cap P_{18}$ and $C=A\cap R$.

\begin{lemma}[witness facts]\label{lem:witness}
Let $w,z,h,r$ be positive integers.
\begin{enumerate}[label=(\roman*)]
\item If $z\equiv7\pmod{25}$, then $25\mid wz+1$ iff $w\equiv7\pmod{25}$. If $z\equiv18\pmod{25}$, then
  $25\mid wz+1$ iff $w\equiv18\pmod{25}$.
\item If $h\equiv\pm7\pmod{25}$ and $r\not\equiv\pm7\pmod{25}$, then $25\nmid hr+1$.
\item $4\mid wz+1$ iff $wz\equiv3\pmod 4$, i.e.\ iff $\{w\bmod4,\,z\bmod4\}=\{1,3\}$.
\end{enumerate}
\end{lemma}
\begin{proof}
(i) $7^{-1}\equiv18$ and $18^{-1}\equiv7\pmod{25}$, so $7w\equiv-1$ iff $w\equiv-18\equiv7$, and similarly for
18. (ii) $h^2\equiv-1$ gives $h^{-1}\equiv-h$, so $25\mid hr+1$ iff $r\equiv-h^{-1}\equiv h\pmod{25}$, which is
excluded. (iii) This is immediate.
\end{proof}

Consequently, if $A$ is a clique, then every element of $Y\cup C$ is a \emph{witness} for $X$: $X$ lies in
its neighbourhood, and for $f\in Y\cup C$ and $z\in X$, $fz+1$ must be divisible by $p^2$ for some prime
$p\ne5$. Likewise $X\cup C$ are witnesses for $Y$, and $X\cup Y$ are witnesses for $C$. Whether $p=2$ is
available is decided by residues mod~4.

\section{Small \texorpdfstring{$N$}{N}}\label{sec:small}

For every $N\le10^4$ we computed $\omega(G_N)$ exactly with a branch-and-bound clique algorithm using greedy
colouring bounds on bitsets (\texttt{maxclique\_small.c}). In all cases $\omega(G_N)=M(N)$. An independent
implementation in Python (Bron--Kerbosch with pivoting and naive squarefree tests, \texttt{bk\_check.py})
confirms this for all $N\le5000$.

\begin{remark}
$G_N$ is not always $M(N)$-colourable \emph{monotonically}. For instance, at $N=43$ the vertices are
$7,18,32,38,41,43$. Here $18\sim43$, $7\sim32$ and $43\sim32$; a 2-colouring exists ($\{7,43,38\}$,
$\{32,18,41\}$), but it cannot be extended consistently from smaller $N$. This is why the certificate below
starts at $N_0=9982$.
\end{remark}

\section{The colouring certificate for \texorpdfstring{$9982\le N\le10^9$}{9982 <= N <= 10^9}}\label{sec:cert}

If the vertex set of $G_N$ is covered by $M(N)$ independent sets, then $\omega(G_N)\le M(N)$. We construct such
covers for all $N$ in the range at once.

\begin{proposition}\label{prop:cert}
Let $N_0=9982$ and $N_1=10^9$. There is a sequence of operations, recorded in a log file, of three types:
\emph{create anchor $x$} ($x\equiv7\bmod25$); \emph{add $v$ to class $x$}; and \emph{move $u$ to class $x$}.
It has the following properties.
\begin{enumerate}[label=(\alph*)]
\item All anchors $x\le N_0$ are created first. Anchors $x>N_0$ are created in increasing order, interleaved
  with the adds.
\item Every vertex $v\le N_1$ of $G_{N_1}$ with $v\not\equiv7\pmod{25}$ is added exactly once. The adds occur
  in increasing order of $v$.
\item Every add or move into class $x$ happens at a time $t$, namely the vertex currently being inserted, with
  $x\le\max(t,N_0)$. The moved or added vertex is non-adjacent to $x$ and to every current member of
  class $x$.
\end{enumerate}
\end{proposition}

\begin{corollary}
For every $N_0\le N\le N_1$, $\omega(G_N)\le M(N)$.
\end{corollary}
\begin{proof}
Consider the state immediately after all vertices $\le N$ have been inserted. By (c), each class is an
independent set, since classes only ever receive vertices non-adjacent to all current members, and removals
preserve independence. By (b), every vertex $\le N$ lies in some class, and that class's anchor is
$\le\max(N,N_0)=N$. Hence $V(G_N)$ is covered by the at most $M(N)$ classes with anchor $\le N$.
\end{proof}

The log is produced by \texttt{dyn\_colour.c}. Each vertex is placed into a random compatible class; if none
is found, one conflicting member is moved to another class (a depth-one augmenting step). Elements of $P_{18}$
are pairwise adjacent and roughly as numerous as anchors, so at almost every time step exactly one class
lacks a $P_{18}$ element; the augmenting moves handle this. For $N_1=10^9$ the run inserts $105{,}158{,}686$
vertices with $13{,}806{,}588$ moves.

The log is checked by \texttt{verify\_log.c}, which was written independently of the generator and uses
different algorithms. It recomputes the vertex set using $\sqrt{-1}\equiv g^{q(q-1)/4}\pmod{q^2}$ (where $g$ is
a non-residue mod $q$), rather than Hensel lifting, and checks every root. It replays every operation and
checks (a)--(c). It decides squarefreeness of every one of the $121{,}112{,}532$ products $ab+1$ claimed
non-adjacent by trial division by all primes up to $(ab+1)^{1/3}$, followed by an exact perfect-square test
of the cofactor. If no prime $p\le n^{1/3}$ divides $n$, then $n$ has at most two prime factors, so it is
squarefree unless it is the square of a prime. The verifier rejects deliberately corrupted logs. All pairs
were additionally re-checked with PARI/GP's \texttt{uissquarefree} (Section~\ref{sec:repro}).

\section{Large \texorpdfstring{$N$}{N}: an explicit argument}\label{sec:analytic}

Throughout, $p,q$ denote primes, $\lambda(w)=\#\{u\bmod w:\ u^2\equiv1\pmod w\}$, and $P_1=P_2=Q=1000$.

\begin{lemma}[large prime squares]\label{lem:LP}
For $1\le w\le N$ and $V\ge1$,
\[
\#\{z\le N:\ p^2\mid wz+1\ \text{for some } p>V\}\ \le\ \lambda(w)\Big(\frac{N^2+1}{V^2}+\frac{2(N+1)}{V}+1\Big).
\]
\end{lemma}
\begin{proof}
If $p^2\mid wz+1$ with $p>V$, write $wz+1=mp^2$. Then $1\le m\le K:=(wN+1)/V^2$ and $mp^2\equiv1\pmod w$. The
pair $(m,p)$ determines $z$. For fixed $m$, if $u_0$ satisfies $mu_0^2\equiv1\pmod w$ then $u_0$ is a unit, and
the solutions mod $w$ are exactly $u_0s$ with $s^2\equiv1$. So the integers $u\le U_m:=\sqrt{(wN+1)/m}$ with
$mu^2\equiv1$ number at most $\lambda(w)(U_m/w+1)$. Summing over $m\le K$ and using
$\sum_{m\le K}m^{-1/2}\le2\sqrt K$ gives at most
$\lambda(w)\big(K+2\sqrt{K(wN+1)}/w\big)=\lambda(w)\big(K+2(wN+1)/(wV)\big)$, which is at most the stated bound
since $w\le N$.
\end{proof}

Writing $w=2^e w'$ with $w'$ odd, $\lambda(w)=c_e2^{\omega(w')}$ with $c_e=1,1,2,4$ for $e=0,1,2,\ge3$. Hence
$\lambda_{\max}(N):=\max_{w\le N}\lambda(w)$ is attained at $2^e$ times a product of initial odd primes, and is
computed exactly.

\begin{lemma}[one witness, $P_7/P_{18}$ target]\label{lem:XY}
Let $Z$ be a residue class modulo $100$ contained in $P_7$ (or in $P_{18}$), and let $w\le N$ satisfy
$25\nmid wz+1$ and $4\nmid wz+1$ for $z\in Z$. Let $\delta=1-\prod_{p\in\{3,7,11,13\}}(1-p^{-2})$. Then for
every $V>P_1$,
\[
\#\{z\in Z\cap[1,N]:\ wz+1\text{ non-squarefree}\}\le \tfrac{N}{100}\Big(\delta+\!\!\sum_{13<p\le P_1}\!\!p^{-2}\Big)
+15+\pi(P_1)-\pi(13) + E_X,
\]
where $E_X$ is the number of $z$ in the whole class $z\equiv c\pmod{25}$ for which $p^2\mid wz+1$ for some
$p>P_1$. Moreover
\[
E_X\le \sum_{P_1<p\le V}\Big(\frac{N}{25p^2}+1\Big)+\lambda(w)\Big(\frac{N^2+1}{V^2}+\frac{2(N+1)}{V}+1\Big).
\]
\end{lemma}
\begin{proof}
The primes $2$ and $5$ are excluded by hypothesis. For $p\in S=\{3,7,11,13\}$ with $p\nmid w$, the condition
$p^2\mid wz+1$ is a single residue class mod $p^2$. By inclusion--exclusion over nonempty subsets $D\subseteq S$,
each term counts a single class mod $100\prod_{p\in D}p^2$ in $[1,N]$, and so differs from its density by at
most $1$. This gives the first term, with error at most $2^4-1=15$; dropping primes dividing $w$ only
decreases the density. For $13<p\le P_1$ use the union bound, each class contributing at most
$N/(100p^2)+1$. For $p>P_1$, bound over the whole residue class mod $25$: at most $N/(25p^2)+1$ for
$P_1<p\le V$, and Lemma~\ref{lem:LP} for $p>V$.
\end{proof}

The quantity $E_X$ depends only on $w$ and the class mod~$25$, so when a single witness bounds several cells
of the same part, $E_X$ is paid once. If $(w\bmod4)(z\bmod4)\equiv3$ the cell is \emph{free}: we only use the
trivial bound $\lfloor N/100\rfloor+1$.

\begin{lemma}[the large-$q$ part of $R$]\label{lem:TR}
Let $L=\{n\le N: q^2\mid n^2+1 \text{ for some } q>Q\}$. For every $V'>Q$,
\[
|L|\le T_R:=\sum_{\substack{Q<q\le V'\\ q\equiv1(4)}}\Big(2\Big\lfloor\frac{N}{q^2}\Big\rfloor+2\Big)
+\frac{N^2+1}{V'^2}\Big(\frac12+\frac{\log(2N+1)}{2\log(1+\sqrt2)}\Big).
\]
\end{lemma}
\begin{proof}
For $q\le V'$, $S_q\cap[1,N]$ consists of two progressions. If $q>V'$ and $q^2\mid n^2+1$, write
$n^2+1=kq^2$ with $k\le(N^2+1)/V'^2$. Here $k\ge2$ is not a square, since $n^2+1=(aq)^2$ is impossible for
$n\ge1$. Positive solutions of $x^2-ky^2=-1$ are units of norm $-1$ in $\mathbb Z[\sqrt k]$, so they are of the
form $\varepsilon^{2j+1}$, where $\varepsilon\ge1+\sqrt2$ is the fundamental unit of that order. Since
$y\sqrt k=\sqrt{x^2+1}\le x+1$, we get $\varepsilon^{2j+1}\le 2x+1\le2N+1$. So for each $k$ there are at most
$\frac12+\log(2N+1)/(2\log(1+\sqrt2))$ values of $n$.
\end{proof}

\begin{definition}
An element $h\in P_7\cup P_{18}$ is \emph{special} if $q^2\mid h^2+1$ for some prime $q\ge13$.
\end{definition}

\begin{lemma}[one witness, $R$ target]\label{lem:C}
Let $h\le N$ with $h\in P_7\cup P_{18}$, let $\tau\in\mathbb Z/4$ with $(h\bmod 4)\tau\not\equiv3$, and let
$V>P_2\ge Q$. The number of $n\in R\setminus L$ with $n\equiv\tau\pmod 4$ and $hn+1$ non-squarefree is at most
\[
\sum_{\substack{13\le q\le Q\\ q\equiv1(4)}}2\Big[\frac{N}{4q^2}\Big(\delta+\!\!\sum_{13<p\le P_2}\!\!p^{-2}\Big)+15+\pi(P_2)-\pi(13)\Big]
+E_C+\mathrm{SP},
\]
where $E_C\le\sum_{P_2<p\le V}(N/p^2+1)+\lambda(h)\big((N^2+1)/V^2+2(N+1)/V+1\big)$. Here $\mathrm{SP}=0$ if $h$
is not special. In general $\mathrm{SP}\le\sum(N/(4q^2)+1)$, summed over the smallest primes
$q\ge13$, $q\equiv1(4)$, whose squares have product at most $(N^2+1)/25$.
\end{lemma}
\begin{proof}
Every $n\in R\setminus L$ lies in one of the classes $\pm i_q\bmod q^2$ with $13\le q\le Q$. Fix such a class.
By Lemma~\ref{lem:witness}, $5$ and $2$ cannot help. For $p\in\{3,7,11,13\}$ and $13<p\le P_2$ with $p\ne q$,
argue as in Lemma~\ref{lem:XY} inside the progression $n\equiv\pm i_q\ (q^2)$, $n\equiv\tau\ (4)$. The prime
$p=q$ is special: on the class, $hn+1$ is constant modulo $q^2$, and $q^2\mid hn+1$ forces
$n\equiv-h^{-1}$, which is compatible with $n^2\equiv-1$ iff $h^2\equiv-1\pmod{q^2}$, i.e.\ iff $h$ is special
at $q$. In that case one whole class $n\equiv h\pmod{q^2}$ is compatible, contributing at most
$N/(4q^2)+1$. The set of $q$ at which $h$ is special satisfies $25\prod q^2\mid h^2+1\le N^2+1$, so its
contribution is at most that of the smallest admissible primes. Primes $p>P_2$ are handled globally, over all
$n\le N$, by Lemma~\ref{lem:LP}.
\end{proof}

\subsection{Patterns}
Let $A\subseteq[1,N]$ be a clique with $A\not\subseteq P_7$ and $A\not\subseteq P_{18}$; otherwise
$|A|\le M(N)$ trivially. Split each of $X$ and $Y$ into four \emph{cells} by residue mod~4. Record, for each
cell, whether it is empty, nonempty with all elements special, or contains a non-special element. For $C$,
record which residues mod~4 occur. This gives $3^4\cdot3^4\cdot2^4$ possible \emph{patterns}.

Given the pattern, we bound each part using at most two witness residues from the other parts. We use any
element of $Y\cup C$ for $X$ and of $X\cup C$ for $Y$. For $C$ we use non-special elements of $X\cup Y$, or
special ones at the cost of SP. Each cell is bounded by the minimum of the applicable bounds:
\begin{itemize}
\item the trivial bound;
\item Lemma~\ref{lem:XY} or~\ref{lem:C} for the chosen witness;
\item for a special-only cell of $X$ or $Y$, the bound $\sum_{13\le q\le Q}2(\lfloor N/(100q^2)\rfloor+1)$,
  since such elements are either in some $S_q$ with $q\le Q$ or lie in $L$.
\end{itemize}
Each witness's $E_X$ or $E_C$ is added once per part. Finally we add $T_R$ once, for $A\cap L$. Elements of
$L$ may be counted twice; this only weakens the bound. Let $\mathrm{UB}(N)$ be the maximum of the resulting
quantity over all patterns. Then $|A|\le\mathrm{UB}(N)$.

\begin{proposition}\label{prop:grid}
There are $5\cdot10^8=N^{(0)}<N^{(1)}<\dots<N^{(k)}=10^{35}$ with
$\mathrm{UB}(N^{(i+1)})<M(N^{(i)})+1$ for all $i$. The maximum of
$\mathrm{UB}(N^{(i+1)})/(M(N^{(i)})+1)$ is $0.9960$, attained on the first interval.
\end{proposition}

For $N^{(i)}\le N\le N^{(i+1)}$, a clique in $G_N$ is also a clique in $G_{N^{(i+1)}}$, so
$|A|\le\mathrm{UB}(N^{(i+1)})<M(N^{(i)})+1\le M(N)+1$. Hence Theorem~\ref{thm:main} holds on
$[5\cdot10^8,10^{35}]$.

\begin{proposition}[tail]\label{prop:tail}
For $N\ge10^{35}$, $\mathrm{UB}(N)<M(N)+1$.
\end{proposition}
\begin{proof}[Sketch]
Take $V=V'=N^{2/3}$. Every term of $\mathrm{UB}(N)$ is either linear in $N$ or sublinear. The linear part is at
most $U_{\rm lin}N$, where $U_{\rm lin}$ is the pattern maximum of the linear coefficients; infinite prime sums
are bounded by exact partial sums up to $3\cdot10^8$ plus $\sum_{n>3\cdot10^8,\,n\ \text{odd}}n^{-2}$. The
sublinear part $S(N)$ comes from at most six witness applications, $T_R$, and the bounded inclusion--exclusion
and rounding constants. It is bounded using $\pi(x)<1.25506\,x/\log x$ \cite{RS62} and
$\lambda(w)\le4\cdot2^{\omega(w)}\le4\cdot2^{1.38402\log N/\log\log N}$ \cite{Ro83}. Each resulting term of
$S(N)/N$ is non-increasing for $N\ge10^8$. Since $M(N)+1>N/25$, it suffices that
$25(U_{\rm lin}+S(10^{35})/10^{35})<1$. In interval arithmetic this quantity is at most $0.7484$, with
$25U_{\rm lin}\le0.6907$.
\end{proof}

The computations behind Propositions~\ref{prop:grid} and~\ref{prop:tail} are in
\texttt{analytic\_rig.py}. All real quantities are evaluated as Arb balls \cite{Arb}, and each is replaced by
the exact rational upper endpoint of its ball; prime counts and floors are exact integers. The maximisation
over patterns is then carried out in exact rational arithmetic. The limiting value $\mathrm{UB}(N)/M(N)\to0.691$
comes from cliques contained in $R$.

\section{Verification and reproducibility}\label{sec:repro}
\begin{itemize}
\item \textbf{Small $N$:} \texttt{maxclique\_small 10000} and \texttt{bk\_check.py 5000}: no mismatches.
\item \textbf{Certificate:} \texttt{dyn\_colour 1000000000 9982 1 log.bin} (about one hour) and
  \texttt{verify\_log 1000000000 9982 log.bin}: VERIFIED, $121{,}112{,}532$ pairs, $0$ failures.
\item \textbf{PARI/GP cross-checks:} PARI's \texttt{uissquarefree} confirms that all
  $121{,}112{,}532$ certificate products $ab+1$ are squarefree. Recomputing the vertex set up to $10^9$ as
  $\{n:\ \texttt{uissquarefree}(n^2+1)=0\}$ gives $105{,}158{,}686$ vertices, with the same sum and the same
  sum of squares (mod $2^{64}$) as the verifier's vertex set.
\item \textbf{Analytic range:} \texttt{analytic\_rig.py grid 5e8 1e35} reports 181 intervals, all
  satisfied, with worst ratio $0.996034$. \texttt{analytic\_rig.py tail 1e35} reports $0.748362<1$. The grid's
  top endpoint and the tail's starting point are the same integer,
  $\lfloor\mathrm{float}(10^{35})\rfloor\approx10^{35}$.
\item \textbf{Independent review:} an independent AI re-implementation of Section~\ref{sec:analytic},
  written without access to our code, found $\mathrm{UB}/M\le0.763$ on $[5\cdot10^8,10^{13}]$. Its lemma
  spot-checks against brute force found no violations.
\end{itemize}

\subsection*{Acknowledgements} We thank M.~Sawhney, W.~van Doorn, D.~Weisenberg, B.~Alexeev, S.~Cambie and
N.~Sothanaphan, whose results and comments on the problem's forum thread this work builds on, and T.~Bloom
for maintaining \texttt{erdosproblems.com}.

\begin{thebibliography}{9}
\bibitem{Er92b} P.~Erd\H{o}s, \emph{Some of my favourite problems in various branches of combinatorics},
  Matematiche (Catania) \textbf{47} (1992), 231--240.
\bibitem{EP848} T.~F.~Bloom, \emph{Erd\H{o}s Problem \#848}, \url{https://www.erdosproblems.com/848}.
\bibitem{Sa25} M.~Sawhney, note on Erd\H{o}s Problem \#848 (2025), linked from \cite{EP848}.
\bibitem{So26} N.~Sothanaphan et al., explicit threshold for Erd\H{o}s Problem \#848 (2026), linked from the
  forum thread of \cite{EP848}.
\bibitem{RS62} J.~B.~Rosser and L.~Schoenfeld, \emph{Approximate formulas for some functions of prime numbers},
  Illinois J.\ Math.\ \textbf{6} (1962), 64--94.
\bibitem{Ro83} G.~Robin, \emph{Estimation de la fonction de Tchebychef $\theta$ sur le $k$-i\`eme nombre premier
  et grandes valeurs de la fonction $\omega(n)$}, Acta Arith.\ \textbf{42} (1983), 367--389.
\bibitem{Arb} F.~Johansson, \emph{Arb: efficient arbitrary-precision midpoint-radius interval arithmetic},
  IEEE Trans.\ Comput.\ \textbf{66} (2017), 1281--1292.
\end{thebibliography}

\end{document}
